% META: {"id":"1SPE-DERLOCAL-CO-011","chapitre":"1SPE-DERIVATION-LOCAL","type_objet":"corrige","exercice_id":"1SPE-DERLOCAL-EX-011","status":"approved","origin":"generated","status_history":["generated","audited","accepted","approved"],"release_acceptance":"RELEASE_OWNER_BATCH_ACCEPTANCE","acceptance_closure_digest":"sha256:9b3ccf9a81c5520fbb7e03b7b2d3e7bf2057a02834b6d908bf3be5f49f3b553f","acceptance_receipt_digest":"sha256:4fad86f0282e0fa4e0419cca1ad6379ae7af5a280c04a4a90880f90a1c044242"}
% BEGIN-VERIFY
% from sympy import symbols, simplify
% a, h = symbols('a h')
% taux = simplify(((a+h)**2 - a**2)/h)
% assert taux == 2*a + h
% END-VERIFY
\begin{corrige}{1SPE-DERLOCAL-EX-011}
\textbf{1.} $(a+h)^2=a^2+2ah+h^2$. Donc $f(a+h)-f(a)=a^2+2ah+h^2-a^2=2ah+h^2$.

\textbf{2.} $\dfrac{f(a+h)-f(a)}{h}=\dfrac{2ah+h^2}{h}=2a+h$.

\textbf{3.} Lorsque $h$ se rapproche de $0$, $2a+h$ se rapproche de $2a$. Donc $f'(a)=2a$ pour tout réel $a$.

\textbf{4.}
\begin{itemize}
  \item $f'(0)=0$ : la tangente au sommet de la parabole est horizontale.
  \item $f'(-1)=-2$ : la tangente descend.
  \item $f'(5)=10$ : la tangente monte.
\end{itemize}

\textbf{5.} $f'(a)=0$ équivaut à $2a=0$, soit $a=0$. La tangente à la parabole au point d'abscisse $0$ est horizontale : c'est le sommet de la parabole, où la courbe change de sens de variation.
\end{corrige}
